\section{Généalogie complète de la projection exceptionnelle : de Paley et
Hadamard aux systèmes de Witt, Leech et Moonshine}
\label{sec:genealogia-excepcional-completa-autonoma}
\index{double projection!généalogie exceptionnelle}
\index{Golay!ternaire et binaire}
\index{Steiner!systèmes de Witt}
\index{Mathieu!groupes M12 et M24}
\index{Moonshine!généalogie}

L'évaluation exceptionnelle part du même état enrichi que
l'évaluation archimédienne. La première conserve l'incidence, l'orientation,
la retenue et la frontière ; la seconde publie la valeur réelle de la section. Cette
bifurcation fait coexister, à partir d'un antécédent commun, une généalogie des
valeurs et une généalogie des symétries.

Deux transformations linéaires remplissent des fonctions successives et typées. La
transformée de Hadamard organise l'état dodécaphasique signé et ses
secteurs cyclotomiques. La matrice de conférence de Paley fixe ensuite le
calibre projectif de la branche exceptionnelle. Sa réduction ternaire
\[
 C_P C_P^{\mathsf T}=5I_6,
 \qquad
 A_W=C_P\bmod3,
 \qquad
 A_W^2=-I_6
\]
produit l'endomorphisme quadratique interne et le code graphique
\[
 C_W=\{(q,qA_W):q\in\mathbb F_3^6\}
      =[12,6,6]_3.
\]
L'énumérateur
\[
 W_{C_W}(y)=1+264y^6+440y^9+24y^{12}
\]
exprime la métrique de Hamming de cette réalisation. Les 264 vecteurs de poids
six, identifiés par paires opposées, déterminent 132 hexades et construisent
le premier système de Witt,
\[
 W_{12}=S(5,6,12).
\]
Les 495 étoiles de Witt fournissent les involutions d'incidence dont
l'action engendrée a pour image \(M_{12}\).

La seconde structure de Steiner est obtenue dans le code binaire étendu
de Golay \(\mathcal G_{24}\). Les supports de poids huit forment
\[
 W_{24}=S(5,8,24),
 \qquad \operatorname{Aut}(W_{24})\cong M_{24}.
\]
Une dodécade \(D\) étant fixée, les intersections des octades avec \(D\) en
six points reconstruisent \(W_{12}\) ; chaque hexade admet un relèvement unique
à une octade. La paire de systèmes
\[
 S(5,6,12)\longrightarrow S(5,8,24)
\]
est donc un prolongement d'incidence démontré au sein du Golay
binaire et conserve tant la dodécade que son alignement.

Du code ternaire commun naît simultanément la branche réticulaire. Le
recollement coordonnée par coordonnée sur \(A_2^{12}\) construit le réseau de
Niemeier
\[
 C_W\longrightarrow N(A_2^{12}).
\]
Le vecteur radial sélectionné définit un voisin d'indice trois pair,
unimodulaire, de rang vingt-quatre et sans racines ; la classification l'identifie
au réseau de Leech :
\[
 N(A_2^{12})\xrightarrow{\;3\text{-voisin}\;}\Lambda_{24}.
\]
Sur cette sortie est construite l'algèbre d'opérateurs de vertex du réseau
\(V_{\Lambda_{24}}\). L'involution canonique de signe et son module tordu
produisent l'orbifold
\[
 V^\natural=V_{\Lambda_{24}}^+
 \oplus V_{\Lambda_{24}}^{T,+},
 \qquad
 \operatorname{Aut}(V^\natural)\cong\mathbb M.
\]
Son caractère gradué est le Hauptmodul normalisé
\[
 \operatorname{Tr}_{V^\natural}q^{L_0-1}
 =J(\tau)=j(\tau)-744
 =q^{-1}+196884q+21493760q^2+\cdots,
\]
et la théorie du Moonshine associe à chaque classe de conjugaison
\([g]\subset\mathbb M\) la série de McKay--Thompson
\(T_g(\tau)=\operatorname{Tr}_{V^\natural}
(g\,q^{L_0-1})\), dotée des propriétés modulaires et de réplicabilité
correspondantes.

L'antécédent commun des deux prolongements conserve la chaîne complète
\[
 w_6\longrightarrow w_{30}\longrightarrow R_{36}\longrightarrow G_9
 \longrightarrow x_{\rm term}\longrightarrow U_{12}^{\rm sgn}
 \longrightarrow K\longrightarrow B_K.
\]
À partir de cet état terminal, il faut distinguer la voie d'évaluation et la voie
du calibre. L'état produit les mots d'incidence ; la matrice de Paley
définit, par réduction ternaire, l'opérateur qui les évalue. La généalogie
exceptionnelle complète s'organise donc comme suit :
\[
\begin{aligned}
C_P&\longmapsto A_W=C_P\bmod3,\\
x\in X_\infty^{\rm cal}
&\longmapsto \varepsilon_\infty(x)=(w_j)_j
 \longmapsto\bigl(\Gamma_{A_j}(w_j)\bigr)_j
 \in\prod_j C_W(A_j),\\
C_W
&\longrightarrow W_{12}\xrightarrow{\ (D,\ell)\ }W_{24},\\
W_{12}&\xrightarrow{\ \operatorname{Aut}\ }M_{12},
\qquad W_{24}\xrightarrow{\ \operatorname{Aut}\ }M_{24},\\
C_W
&\longrightarrow N(A_2^{12})\longrightarrow\Lambda_{24}
 \longrightarrow V_{\Lambda_{24}}\longrightarrow V^\natural
 \longrightarrow\mathbb M\longrightarrow\text{Moonshine}.
\end{aligned}
\]
Le poids et la distance de Hamming appartiennent à la métrique des codes ;
Hadamard appartient à la transformation dodécaphasique ; Paley fixe le calibre
de conférence. Cette séparation des fonctions permet de suivre chaque flèche par
son opérateur, son domaine et son certificat, tout en conservant l'unité
généalogique de la double projection.
